\documentclass[12pt]{article}
\usepackage[utf8]{inputenc}
\usepackage{mathptmx} % Set Times New Roman font
\usepackage{graphicx} % Required for including images
\usepackage{fancyhdr} % Required for custom headers and footers
\usepackage{booktabs}
\usepackage[margin=1in, headheight=33pt]{geometry} % Set standard 1-inch margins
\usepackage[table]{xcolor}
\usepackage{colortbl}
\definecolor{lightgray}{gray}{0.9}
\usepackage{pdfpages}
\usepackage{hyperref}
\usepackage{xcolor}
\usepackage{caption}
\usepackage{float} % <-- for [H] placement

\hypersetup{
    colorlinks=true,
    linkcolor=black,
    filecolor=black,
    urlcolor=black,
    citecolor=black
}

% Header and Footer Settings
\pagestyle{fancy}
\fancyhf{}
\fancyhead[L]{\includegraphics[height=1cm]{../png/uga-logo.png}}
\fancyhead[R]{Bohdan Leniuk}
\renewcommand{\headrulewidth}{0.4pt}
\fancyfoot[C]{\thepage}
\renewcommand{\footrulewidth}{0pt}
\setlength{\parindent}{0.5in}

\captionsetup[table]{skip=5pt}

\fancypagestyle{firstpage}{
  \fancyhf{}
  \fancyhead[L]{\includegraphics[height=1cm]{../png/uga-logo.png}}
  \fancyhead[R]{Bohdan Leniuk}
  \renewcommand{\headrulewidth}{0.4pt}
  \fancyfoot{}
}

\begin{document}
\thispagestyle{firstpage}

\begin{center}
    {\large \textbf{Teaching Statement}} \\[0.5em]
    Bohdan Leniuk
\end{center}

\vspace{0.5em}

As an instructor of economics, I believe that my primary responsibility is to cultivate genuine intellectual curiosity in my students. This is not to say that performance is unimportant, but that it is a byproduct for any student who is genuinely curious about the content they are learning. Economics is, after all, a discipline that provides a framework for understanding the world around us -- one that is not easily learned by memorizing formulas or definitions. By fostering that curiosity, students become self-motivated to learn not merely the material at hand, but the mode of thinking required for success in courses to come and beyond.  

This is not a simple task. Within a single classroom, there are both those who come in already curious, eager to learn more about a passion they already possess, and those who come in not yet so, eager only to fulfill a requirement for graduation. Optimal treatment for such students should, in my experience, differ because \textit{developing} a sense of curiosity is not the same as \textit{sustaining} one. A successful educator can both identify which students are which and can address their heterogeneous needs within the classroom. 

Students who enter without a prior connection to the subject respond best when instruction meets them where their existing interests lie. Discussions within the classroom, examples, and even practice problems should reflect the interests of the students in question. Not only does this give students a reason to engage, but, perhaps more importantly, it creates a link between economics and something they are already connected with. Once this link is established, a student who once viewed economics as a box to check can begin to see it as a lens through which the world makes more sense. The curiosity that follows tends to be genuine and, importantly, self-sustaining. 

One strategy I have found particularly effective in this regard is to construct examples collaboratively with students during lecture. Rather than arriving with a fixed scenario prepared, I allow the class to choose the context -- the market, the goods being traded, the agents involved. This accomplishes two things: it ensures the example is relevant to the students in the room, and it gives each student a degree of ownership over the material we are working through together. A particularly clear instance of this arose when teaching consumer demand to a section that included several art history majors. Taking their lead, we examined the market for Banksy's work in the aftermath of his 2018 auction stunt, in which his \textit{Girl with Balloon} partially shredded itself moments after selling at Sotheby's. The incident generated substantial public exposure, and the market value of his work rose noticeably in the weeks that followed. The physical materials of the piece had not changed; nor had the work simply appreciated over time. What changed was exposure -- which expanded the pool of potential buyers, shifted the demand curve outward, and drove up the equilibrium price. The example mapped cleanly onto the theory, and the students for whom it was designed were visibly engaged. Worth noting, however, is that I came to this example through them -- the Banksy incident was not something I had previously encountered, and learning of it in the course of preparing for their classroom was itself instructive. I have come to believe that this dynamic is not incidental: when students observe that I am genuinely curious about their world, the expectation that they engage with mine becomes considerably easier to fulfill. 

On the other hand, there are always a select few who come into the class with an already existing passion. Such students do not need help developing a connection to economics. On the contrary, I find that time spent on such efforts serve only to bore the student and, in the worst case, erode any pre-existing interest. These students need to be challenged. Discussions, examples, and practice problems should be sufficiently nuanced such that they deepen their understanding of the underlying theory. The risk, of course, is that in designing instruction around depth and nuance, one can easily neglect those students who are still searching for their initial connection to the discipline altogether.

Therefore, if one aims to serve both types of students effectively, it is imperative to identify early which students need which approach. In my own experience, I have attempted to address this challenge with a pre-semester survey. Within it, students were asked about their relationship to economics, their reasons for enrolling, and the fields they study outside of my own classroom. I found this to be initially beneficial, as it provided me with an early and practical map of my students -- who needed a bridge connecting them to the material, and who needed to be pushed further within it. Well-intentioned as this was, I do not believe it was perfect. As the semester carried on, I found that my students grew rapidly across this dimension. Students who began the semester disengaged became invested over time, and as they did, the instruction they required shifted accordingly. This taught me an ironic lesson. While it is my duty to teach, doing so effectively requires that I study my students both before and throughout the course. Actively doing so, and adjusting the classroom along with it, is what is required of a teacher to achieve the goal of genuine curiosity.  


\end{document}
